% ===== PRÉAMBULE CANONIQUE — à coller VERBATIM en tête de CHAQUE .tex =====
\documentclass[11pt,a4paper]{article}
\usepackage[utf8]{inputenc}\usepackage[T1]{fontenc}\usepackage{lmodern}
\usepackage[french]{babel}
\usepackage{amsmath,amssymb,mathtools}\usepackage{icomma}
\usepackage[version=4]{mhchem}          % \ce{...} : équations & formules
\usepackage{siunitx}\sisetup{locale=FR} % \qty{}{}, \si{}, \ang{}
\usepackage{chemfig}                    % \chemfig{...} : structures développées
\usepackage{xcolor}\usepackage{colortbl}
\usepackage{tikz}\usepackage{pgfplots}\pgfplotsset{compat=1.18}
\usepackage[most]{tcolorbox}\usepackage{enumitem}\usepackage{array,booktabs}
\usepackage[a4paper,margin=2.2cm]{geometry}
\usetikzlibrary{arrows.meta,calc,angles,quotes,positioning,patterns,backgrounds,shapes.geometric}
\definecolor{primary}{RGB}{222,49,99}\definecolor{primarydark}{RGB}{150,22,55}
\definecolor{defcol}{RGB}{0,114,178}\definecolor{methcol}{RGB}{52,92,148}
\definecolor{excol}{RGB}{0,135,90}\definecolor{attncol}{RGB}{200,40,55}
\tcbset{boxbase/.style={enhanced,breakable,boxrule=0.9pt,arc=2.5pt,
  left=3mm,right=3mm,top=2.3mm,bottom=2.3mm,fonttitle=\bfseries,coltitle=white}}
% --- boîtes TUTORIEL (noms EXACTS, ne pas renommer) ---
\newtcolorbox{cle}[1][]{boxbase,colback=defcol!6,colframe=defcol,
  title={\textsc{À retenir}\ifx\relax#1\relax\relax\else\textmd{ : \itshape #1}\fi}}
\newtcolorbox{methode}[1][]{boxbase,colback=methcol!7,colframe=methcol,sharp corners,
  title={\textsc{Méthode}\ifx\relax#1\relax\relax\else\textmd{ --- \itshape #1}\fi}}
\newtcolorbox{exemplebox}[1][]{boxbase,colback=excol!5,colframe=excol,
  title={\textsc{Exemple}\ifx\relax#1\relax\relax\else\textmd{ : \itshape #1}\fi}}
\newtcolorbox{piege}[1][]{boxbase,colback=attncol!6,colframe=attncol,
  title={\textsc{Piège}\ifx\relax#1\relax\relax\else\textmd{ : \itshape #1}\fi}}
\newtcolorbox{remarque}[1][]{boxbase,colback=black!4,colframe=black!45,
  title={\textsc{Remarque}\ifx\relax#1\relax\relax\else\textmd{ : \itshape #1}\fi}}
% --- boîtes EXERCICE / CORRIGÉ (noms EXACTS, auto-comptées) ---
\newtcolorbox[auto counter]{exo}[1][]{boxbase,colback=primary!4,colframe=primary,
  title={\textsc{Exercice}~\thetcbcounter\ifx\relax#1\relax\relax\else\textmd{ --- \itshape #1}\fi}}
\newtcolorbox[auto counter]{corrige}[1][]{boxbase,colback=excol!4,colframe=excol,
  title={\textsc{Corrigé}~\thetcbcounter\ifx\relax#1\relax\relax\else\textmd{ --- \itshape #1}\fi}}
\newcommand{\R}{\mathbb{R}}\newcommand{\N}{\mathbb{N}}\newcommand{\Z}{\mathbb{Z}}\newcommand{\ds}{\displaystyle}
% (un \usepackage{fancyhdr}+en-tête est possible ; il est ignoré côté web)
% =========================================================================
\begin{document}
\sisetup{per-mode=symbol}
\begin{corrige}[l'ambiguïté des entiers]
Un entier terminé par des zéros est ambigu : on ne sait pas si ces zéros résultent d'une mesure. La
notation scientifique lève l'ambiguïté, car la mantisse ne contient que les CS.
\begin{enumerate}[label=\alph*)]
  \item $200$ : $1$, $2$ ou $3$ CS possibles. À $3$ CS : $\boxed{2,00\times10^{2}}$.
  \item $4500$ : $2$, $3$ ou $4$ CS possibles. À $2$ CS : $\boxed{4,5\times10^{3}}$.
  \item $70$ : $1$ ou $2$ CS possibles. À $2$ CS : $\boxed{7,0\times10^{1}}$.
\end{enumerate}
\end{corrige}

\begin{corrige}[lire les CS sur la notation scientifique]
On compte simplement les chiffres de la mantisse ; la puissance de dix n'apporte aucun CS.
\begin{enumerate}[label=\alph*)]
  \item $4,00\times10^{3}$ : $\mathbf{3}$ CS.
  \item $1,0\times10^{-7}$ : $\mathbf{2}$ CS.
  \item $6,022\times10^{23}$ : $\mathbf{4}$ CS.
  \item $2,50\times10^{-2}$ : $\mathbf{3}$ CS.
\end{enumerate}
\end{corrige}

\begin{corrige}[le changement d'unité ne change pas les CS]
Passer d'une unité à l'autre ne fait que déplacer la virgule : le nombre de CS est conservé.
\begin{enumerate}[label=\alph*)]
  \item $\num{1,302e4}\ \si{\milli\gram}$ : $\mathbf{4}$ CS.
  \item $\qty{0,01302}{\kilo\gram}$ : $\mathbf{4}$ CS (les deux zéros de tête ne comptent pas).
  \item $\num{1,302e-5}\ \si{\tonne}$ : $\mathbf{4}$ CS.
\end{enumerate}
\textbf{Conclusion.} Les trois écritures ont $4$ CS, comme $\qty{13,02}{\gram}$ : la précision de la
mesure n'a pas changé, seule l'écriture a changé.
\end{corrige}

\begin{corrige}[mesuré ou exact ?]
Un nombre \emph{compté} ou \emph{défini} est exact (infinité de CS) ; un nombre \emph{lu sur un
instrument} est mesuré (nombre fini de CS).
\begin{enumerate}[label=\alph*)]
  \item \textbf{Exact} (comptage d'atomes) : infinité de CS.
  \item \textbf{Mesuré} : $\qty{24,15}{\milli\litre}$ a $4$ CS.
  \item \textbf{Exact} (définition) : infinité de CS.
  \item \textbf{Mesuré} : $\qty{0,250}{\gram}$ a $3$ CS (le zéro de queue après la virgule compte).
\end{enumerate}
\end{corrige}

\begin{corrige}[comparer la précision]
La précision se lit sur le nombre de CS, indépendamment de l'unité. On compte :
\begin{enumerate}[label=\alph*)]
  \item $\qty{0,04500}{\metre}$ : $4$ CS.
  \item $\qty{4,5}{\centi\metre}$ : $2$ CS.
  \item $\qty{45,0}{\milli\metre}$ : $3$ CS.
\end{enumerate}
Ordre de la moins à la plus précise : $\text{b)}\ (2\ \text{CS}) < \text{c)}\ (3\ \text{CS}) <
\text{a)}\ (4\ \text{CS})$. Bien que ce soit la \emph{même} longueur, c'est l'écriture qui affiche
le plus de CS qui traduit la mesure la plus fine.
\end{corrige}

\end{document}
